\chapter{System Design}

\section{System Architecture}
GeoCrop follows a layered, service-oriented architecture built around one principle: \emph{send
only what is visible}. Model results are stored once as Cloud Optimized GeoTIFFs (COG) in object
storage; the browser requests small map tiles for the current view, which a stateless tile service
cuts from the COG and Nginx caches. Heavy work runs in background workers, and a spatial database
holds the catalogue that tells the interface what exists. The main building blocks are the COG
format~\cite{cog}, TiTiler~\cite{titiler}, PMTiles~\cite{pmtiles}, MapLibre GL~\cite{maplibre} and
PostGIS~\cite{postgis}. Figure~\ref{fig:arch} shows the services
and their connections.

\begin{figure}[htbp]
\centering
\begin{tikzpicture}[every node/.style={font=\small}]
\node[box, minimum width=4.2cm, minimum height=0.7cm, fill=gray!8] (base) at (0,1.35) {Basemap tiles (OSM / MapTiler)};
\node[box, minimum width=11.4cm, minimum height=0.95cm, fill=green!10] (client) at (0,0) {\textbf{Web GIS client}: React + MapLibre GL (react-map-gl)};
\node[box, minimum width=11.4cm, minimum height=0.85cm, fill=yellow!10] (nginx) at (0,-1.6) {\textbf{Nginx}: serves the app, proxies /api, caches /tiles};
\node[box, text width=3.1cm, minimum height=1.35cm] (api) at (-3.95,-3.55) {\textbf{FastAPI}\\auth, catalogue, statistics, uploads, jobs};
\node[box, text width=3.1cm, minimum height=1.35cm] (tiles) at (0,-3.55) {\textbf{Tile service}\\TiTiler / rio-tiler\\COG $\rightarrow$ PNG tiles};
\node[box, text width=3.1cm, minimum height=1.35cm] (worker) at (3.95,-3.55) {\textbf{Celery workers}\\ingest, preprocessing, model registry runs};
\node[dbox, minimum width=2.8cm] (pg) at (-3.95,-7.0) {PostgreSQL + PostGIS\\catalogue, statistics};
\node[dbox, minimum width=2.8cm] (s3) at (0,-7.0) {MinIO (S3)\\COG, PMTiles, uploads};
\node[dbox, minimum width=2.8cm] (redis) at (3.95,-7.0) {Redis\\job broker};
\node[lane, fit=(pg)(s3)(redis), inner sep=7pt, label={[font=\scriptsize\itshape, anchor=north west]north west:Data tier}] (data) {};
\draw[arr] (client) -- (base);
\draw[arr] (client) -- (nginx);
\draw[arr] (nginx.south -| api.north) -- (api.north);
\draw[arr] (nginx.south -| tiles.north) -- (tiles.north);
\draw[darr] (tiles.west) -- node[lbl,above]{layer} (api.east);
\draw[arr] (api.south) -- node[lbl,right]{SQL, enqueue} (api.south |- data.north);
\draw[arr] (tiles.south) -- node[lbl,right]{COG range reads} (tiles.south |- data.north);
\draw[arr] (worker.south) -- node[lbl,right]{jobs, results, statistics} (worker.south |- data.north);
\end{tikzpicture}
\caption{GeoCrop system architecture. The browser also reads PMTiles and uploads imagery directly to object storage using short-lived signed URLs; the tile service looks up layer styles from the API.}
\label{fig:arch}
\end{figure}

\section{Architectural Components}
\begin{table}[htbp]
\centering\small
\caption{Architectural components}
\label{tab:components}
\begin{tabularx}{\textwidth}{L{3.0cm} L{3.4cm} Y}
\toprule
\textbf{Component} & \textbf{Technology} & \textbf{Responsibility}\\
\midrule
Web GIS client & React, Vite, TypeScript, MapLibre GL via react-map-gl, TanStack Query, Zustand & Home, Live Map workspace, Dashboard, Analyze Imagery, Processing, Models, Status, accounts. Builds layers from the catalogue; loads tiles on demand.\\
Edge & Nginx & Serves the built app; proxies the API; caches tiles with immutable headers; resolves services through Docker DNS at request time.\\
API service & FastAPI, Pydantic, SQLAlchemy 2, GeoAlchemy2, PyJWT, bcrypt & REST API v1: health, auth, areas, models, layers, fields, point query, statistics, uploads, runs, jobs.\\
Tile service & TiTiler / rio-tiler (FastAPI) & \code{/tiles/{layer}/{version}/{z}/{x}/{y}.png}; reads only needed COG blocks; applies the server-side style; optional crop filter.\\
Workers & Celery & Ingest (validate, reproject, COG), model runs through the registry, statistics, layer publishing, progress.\\
Model registry & Python package with YAML manifests & Discovers models, resolves dependencies, runs tasks with a common context; built-in cropland, crop type, NDVI index, boundary and yield placeholders.\\
Spatial database & PostgreSQL 16 + PostGIS 3.4, Alembic & Catalogue of users, areas, uploads, images, models, runs, layers, statistics, field regions, jobs.\\
Job broker & Redis 7 & Celery broker and result backend.\\
Object storage & MinIO (S3 API) & Uploads, COG imagery and results, PMTiles; anonymous read only for vector tiles.\\
\bottomrule
\end{tabularx}
\end{table}

\section{Interfaces and Dependencies}
\begin{table}[htbp]
\centering\small
\caption[Main REST interfaces]{Main REST interfaces (prefix \code{/api/v1}; interactive documentation at \code{/api/docs})}
\label{tab:api}
\begin{tabularx}{\textwidth}{L{5.2cm} Y}
\toprule
\textbf{Endpoint} & \textbf{Purpose}\\
\midrule
\code{GET /health} & Status of database, broker, storage, tiles and model registry; configuration notes.\\
\code{POST /auth/signup}, \code{/auth/login}, \code{GET /auth/me} & Accounts and tokens.\\
\code{GET /aois} & Areas visible to the user, with bounds.\\
\code{GET /models} & Registered models, status, inputs, outputs, accuracy and per-class scores.\\
\code{GET /layers?aoi=&kind=&run=} & Layer catalogue: style, legend, bounds, zoom range, tile URL, version.\\
\code{GET /fields?aoi=&q=} & Field regions with crop, area and bounding box.\\
\code{GET /point?lng=&lat=&layers=} & Values and labels of several layers at one location.\\
\code{GET /stats?layer=} & Pre-computed area and share per class.\\
\code{POST /uploads}, \code{/uploads/{id}/complete} & Signed upload URL; start ingestion with models, location and date.\\
\code{POST /runs}, \code{GET /runs}, \code{GET /jobs/{id}} & Re-run models; run history; job progress and steps.\\
\code{GET /tiles/{layer}/{v}/{z}/{x}/{y}.png} & Raster tiles (tile service, cached by Nginx).\\
\bottomrule
\end{tabularx}
\end{table}
External dependencies are Sentinel-2 imagery exported with Google Earth Engine, USDA CDL rasters and
AOI shapefiles, an OpenStreetMap or MapTiler basemap, and the open-source libraries listed in
Section~\ref{sec:external}. Internally, the API depends on the database, broker and storage; the tile
service depends on the API (layer lookup, cached for five minutes) and storage; workers depend on the
broker, database, storage and model artifacts.

\section{Architecture Decisions and Rationale}
\begin{table}[htbp]
\centering\small
\caption{Key architecture decisions}
\label{tab:adr}
\begin{tabularx}{\textwidth}{L{3.6cm} Y L{3.6cm}}
\toprule
\textbf{Decision} & \textbf{Rationale} & \textbf{Alternative rejected}\\
\midrule
Raster as source of truth (COG) & Models predict pixels; a 30~m raster of AOI1 is about 1.6~MB, the same result as polygons 43.3~MB. & Vectorising all output to GeoJSON\\
Tiles via TiTiler & Reads only the COG blocks needed per tile; styles applied on the server; no pre-rendering. & Pre-rendered tile folders; single PNG overlays\\
PMTiles for vectors & One static file read by range requests; no tile server; cacheable. & GeoJSON; dynamic vector tile server\\
MapLibre GL & WebGL rendering of raster and vector tiles; open source. & Leaflet (SVG/Canvas)\\
Versioned, immutable tile URLs & Browsers and Nginx can cache for a year; a new run gets a new URL. & Cache invalidation logic\\
FastAPI & Same stack as TiTiler; async I/O; typed schemas and generated docs. & Django templates (prototype)\\
PostgreSQL + PostGIS & Spatial types and indexes, joins between catalogue and statistics, JSONB for styles and metrics. & SQLite (prototype)\\
Celery + Redis & Long jobs off the request path; retries; progress. & Synchronous processing\\
S3-compatible storage (MinIO) & Large files, range reads, signed uploads, cloud-portable. & Local disk\\
Manifest-based model registry & Models change over the project; new models without changes elsewhere. & Hard-coded model calls\\
Docker Compose & One-command, reproducible stack for the team and the panel. & Manual installation\\
\bottomrule
\end{tabularx}
\end{table}

\section{Design Evolution}\label{sec:evolution}
The proposal-stage prototype was a Django application with Leaflet. It returned the crop-type result
for AOI1 as one GeoJSON response; the browser parsed it and drew 101,435 polygons as SVG. Uploads
were recorded but never processed, field data on the dashboard was hand-written, and the session
cookie held upload history. Table~\ref{tab:evolution} summarises the changes and why they were made.

\begin{table}[htbp]
\centering\small
\caption{Design evolution since the proposal}
\label{tab:evolution}
\begin{tabularx}{\textwidth}{L{2.6cm} L{3.6cm} L{4.0cm} Y}
\toprule
\textbf{Aspect} & \textbf{Proposal prototype} & \textbf{Current design} & \textbf{Reason}\\
\midrule
Result format & GeoJSON polygons (43.3~MB) & COG rasters + PMTiles & Data volume and rendering cost\\
Map client & Leaflet, SVG & MapLibre GL, WebGL tiles & Smooth zoom with many layers\\
Backend & Django views and templates & FastAPI + React SPA & Async API, typed contracts, separation of concerns\\
Database & SQLite & PostgreSQL + PostGIS & Spatial queries, concurrency, catalogue\\
Processing & None (upload recorded only) & Celery jobs with progress & Real analysis of new imagery\\
Models & Offline scripts writing GeoJSON & Registry plug-ins writing rasters & Add models without code changes\\
Dashboard data & Hard-coded sample fields & Pre-computed statistics in PostGIS & Real results\\
Deployment & Single process & Seven containers via Compose & Reproducible, scalable services\\
\bottomrule
\end{tabularx}
\end{table}

\section{Domain Model}
Figure~\ref{fig:domain} shows the main domain classes. An \emph{Area of Interest} owns source images
and model runs; a run executes one registered \emph{Model} and produces one or more \emph{Layers};
categorical layers are summarised by \emph{Zonal statistics}; crop-type runs also produce
\emph{Field regions}. A \emph{Job} tracks the steps of an upload or a run.

\begin{figure}[htbp]
\centering
\resizebox{\textwidth}{!}{%
\begin{tikzpicture}[
  cls/.style={draw=black!70, rectangle split, rectangle split parts=2, fill=green!4, font=\scriptsize, align=left, text width=3.0cm, inner sep=3pt, rectangle split part fill={green!14,green!3}},
  rel/.style={thick, black!60}, m/.style={font=\scriptsize, inner sep=1pt}]
\node[cls] (user) at (0,0) {\textbf{User}\nodepart{second}email, username\\password\_hash, role};
\node[cls] (aoi) at (4.4,0) {\textbf{AreaOfInterest}\nodepart{second}name, geom (Polygon)\\is\_public};
\node[cls] (upload) at (0,-2.6) {\textbf{Upload}\nodepart{second}filename, size\\storage\_key, status};
\node[cls] (img) at (4.4,-2.6) {\textbf{SourceImage}\nodepart{second}kind, captured\_on\\bands, resolution, bounds};
\node[cls] (model) at (8.8,0) {\textbf{Model}\nodepart{second}key, task, version\\status, manifest};
\node[cls] (run) at (8.8,-2.6) {\textbf{ModelRun}\nodepart{second}status, params\\metrics, timings};
\node[cls] (layer) at (8.8,-5.3) {\textbf{Layer}\nodepart{second}kind, format, storage\_key\\style, legend, version};
\node[cls] (stat) at (4.4,-5.3) {\textbf{ZonalStat}\nodepart{second}class, pixel\_count\\area\_ha, mean};
\node[cls] (field) at (13.0,-2.6) {\textbf{FieldRegion}\nodepart{second}geom (Polygon)\\crop, area\_ha};
\node[cls] (job) at (0,-5.3) {\textbf{Job}\nodepart{second}kind, status, progress\\steps, error};
\draw[rel] (user) -- node[m,above]{owns} node[m,pos=0.15,below]{1} node[m,pos=0.85,below]{*} (aoi);
\draw[rel] (user) -- node[m,left]{submits} node[m,pos=0.15,right]{1} node[m,pos=0.85,right]{*} (upload);
\draw[rel] (upload) -- node[m,above]{becomes} node[m,pos=0.15,below]{1} node[m,pos=0.85,below]{0..1} (img);
\draw[rel] (aoi) -- node[m,right]{has} node[m,pos=0.15,left]{1} node[m,pos=0.85,left]{*} (img);
\draw[rel] (model) -- node[m,right]{runs as} node[m,pos=0.15,left]{1} node[m,pos=0.85,left]{*} (run);
\draw[rel] (img) -- node[m,above]{input of} node[m,pos=0.15,below]{1} node[m,pos=0.85,below]{*} (run);
\draw[rel] (run) -- node[m,right]{produces} node[m,pos=0.15,left]{1} node[m,pos=0.85,left]{1..*} (layer);
\draw[rel] (layer) -- node[m,above]{summarised by} node[m,pos=0.15,below]{1} node[m,pos=0.85,below]{*} (stat);
\draw[rel] (run) -- node[m,above]{derives} node[m,pos=0.15,below]{1} node[m,pos=0.85,below]{*} (field);
\draw[rel] (upload) -- node[m,left]{tracked by} node[m,pos=0.15,right]{1} node[m,pos=0.85,right]{*} (job);
\end{tikzpicture}}
\caption{Domain model}
\label{fig:domain}
\end{figure}

\section{Data Design}
\textbf{Relational schema.} The catalogue is a PostgreSQL/PostGIS schema managed with Alembic
(Table~\ref{tab:schema}). Geometries are stored in EPSG:4326 with GiST spatial indexes. Free-form
fields (layer kind, style, legend, model manifest, run metrics) are strings or JSONB, so a new model
output never needs a schema change. Rasters are never stored in the database.

\begin{table}[htbp]
\centering\small
\caption{Database tables}
\label{tab:schema}
\begin{tabularx}{\textwidth}{L{2.4cm} Y}
\toprule
\textbf{Table} & \textbf{Key columns}\\
\midrule
\code{users} & id, email (unique), username (unique), password\_hash, role, created\_at\\
\code{aoi} & id, owner\_id, name, geom (Polygon, GiST), is\_public\\
\code{upload} & id, user\_id, aoi\_id, filename, size\_bytes, storage\_key, status\\
\code{source_image} & id, aoi\_id, upload\_id, kind, captured\_on, storage\_key, bounds (GiST), band\_count, resolution\_m\\
\code{model} & id, key (unique), task, version, name, status, manifest (JSONB), enabled\\
\code{model_run} & id, aoi\_id, model\_id, source\_image\_id, status, params, metrics (JSONB), error, timings\\
\code{map_layer} & id, run\_id, aoi\_id, kind, group, name, format, storage\_key, bounds (GiST), zoom range, nodata, style, legend (JSONB), version, is\_public\\
\code{zonal_stat} & layer\_id, zone\_type, zone\_id, class\_code, class\_name, pixel\_count, area\_ha, mean, min, max\\
\code{field} & id, aoi\_id, run\_id, geom (Polygon, GiST), attributes (JSONB)\\
\code{job} & id, kind, status, progress, message, error, steps (JSONB), upload\_id, user\_id\\
\bottomrule
\end{tabularx}
\end{table}

\textbf{Object storage layout} (bucket \code{geocrop}): \code{raw/{upload}/original.tif} for
uploads; \code{runs/reference/...cog.tif} for imagery and reference layers;
\code{runs/{run}/{kind}-{id}.cog.tif} for model outputs; and \code{vector/{aoi}/{run}/...pmtiles}
for vector layers. Files are written once; a new run creates new keys.

\textbf{Raster conventions.} COGs use 512-pixel internal blocks, DEFLATE compression and overviews
(nearest neighbour for class maps, average for imagery). Categorical rasters use 0 as no-data; the
cropland map uses 1 = cropland and 2 = non-cropland; crop-type codes map to class names stored in the
layer legend.

\section{Process / Behavior Design}
\textbf{Upload to map.} Figure~\ref{fig:seq} shows the interaction from upload to published layers.

\begin{figure}[htbp]
\centering
\begin{tikzpicture}[font=\scriptsize]
\foreach \x/\n/\t in {0/b/Browser, 2.9/a/FastAPI, 5.8/s/MinIO, 8.7/w/Worker, 11.6/d/PostGIS}{
  \node[box, minimum width=2.1cm] (\n) at (\x,0) {\textbf{\t}};
  \draw[dashed, black!40] (\x,-0.35) -- (\x,-10.1);
}
\newcommand{\msg}[5]{\draw[#5] (#1,#3) -- node[lbl,above]{#4} (#2,#3);}
\msg{0}{2.9}{-0.9}{1. POST /uploads}{arr}
\msg{2.9}{11.6}{-1.4}{2. insert upload}{arr}
\msg{2.9}{0}{-1.9}{3. signed PUT URL}{darr}
\msg{0}{5.8}{-2.5}{4. PUT image (direct)}{arr}
\msg{0}{2.9}{-3.1}{5. POST /complete (models)}{arr}
\msg{2.9}{11.6}{-3.6}{6. insert job + steps}{arr}
\msg{2.9}{8.7}{-4.1}{7. enqueue ingest}{arr}
\msg{8.7}{5.8}{-4.7}{8. read original}{arr}
\draw[arr] (8.7,-5.2) -- ++(0.7,0) -- node[lbl,right]{9. validate, reproject, COG} ++(0,-0.45) -- ++(-0.7,0);
\msg{8.7}{5.8}{-6.1}{10. write image COG}{arr}
\msg{8.7}{11.6}{-6.6}{11. AOI, image, layer}{arr}
\draw[arr] (8.7,-7.1) -- ++(0.7,0) -- node[lbl,right]{12. run models in order} ++(0,-0.45) -- ++(-0.7,0);
\msg{8.7}{5.8}{-8.0}{13. result COGs}{arr}
\msg{8.7}{11.6}{-8.5}{14. run, layers, statistics}{arr}
\msg{0}{2.9}{-9.1}{15. GET /jobs/{id} (polling)}{arr}
\msg{2.9}{0}{-9.6}{16. done, steps}{darr}
\end{tikzpicture}
\caption{Sequence of an upload and analysis}
\label{fig:seq}
\end{figure}

\textbf{Job and run states.} A job moves from \emph{queued} to \emph{running} and ends in \emph{done}
or \emph{failed}; each step inside it is \emph{pending}, \emph{running}, \emph{done}, \emph{skipped}
(placeholder model) or \emph{failed} (Figure~\ref{fig:states}).

\begin{figure}[htbp]
\centering
\begin{tikzpicture}[st/.style={draw=black!70, rounded corners=8pt, fill=green!5, minimum width=2.1cm, minimum height=0.75cm, font=\small}]
\node[circle, fill=black, inner sep=2.5pt] (start) at (0,0) {};
\node[st] (q) at (2,0) {queued};
\node[st] (r) at (6,0) {running};
\node[st] (d) at (12,1.1) {done};
\node[st, fill=red!6] (f) at (12,-1.1) {failed};
\draw[arr] (start) -- (q);
\draw[arr] (q) -- node[lbl,above]{worker picks job} (r);
\draw[arr] (r.east) -- ++(0.6,0) |- node[lbl,pos=0.75,above]{all steps done} (d.west);
\draw[arr] (r.east) -- ++(0.6,0) |- node[lbl,pos=0.75,below]{invalid input or error} (f.west);
\draw[arr] (r.south) .. controls +(-0.6,-1.0) and +(0.6,-1.0) .. node[lbl,below]{progress per step} (r.south);
\end{tikzpicture}
\caption{Job state machine}
\label{fig:states}
\end{figure}

\textbf{Model execution.} The registry plans the requested models with a topological sort over the
layer kinds they require and produce (for example, crop type requires cropland). Each model receives
a run context (image path, upstream outputs, scratch folder, artifacts, progress callback); its
output files are converted to COG, uploaded, registered as layers with style and legend, and
summarised as zonal statistics before the next model starts, so the map fills in progressively.

\textbf{Tile request.} The client requests \code{/tiles/{layer}/{version}/{z}/{x}/{y}.png}. Nginx
returns a cached copy when it has one; otherwise the tile service looks up the layer (cached),
reads the needed COG window from storage with nearest-neighbour resampling for class maps, applies
the stored colour map and returns a PNG with a one-year immutable cache header. Tiles outside the
raster return HTTP~204.
